\chapter*{Abstract}
\thispagestyle{empty}
This thesis develops and documents a Pakistan-adapted reinforcement learning workflow for long-horizon crop-management experimentation on top of the CyclesGym environment. The work begins from the original CyclesGym base repository and extends it in five implementation phases: Pakistan crop-calendar alignment, Pakistan price localization, multi-nutrient fertilization support, hierarchical integration of yearly crop planning with within-season fertilization, and standardized reporting for thesis-grade evaluation. The implemented stack uses a Pakistan weather file, Pakistan soil profile, and yearly crop and nutrient price series reconstructed from official sources, while keeping the active experimental crop setup on the current maize-soy configuration that is already wired into the repository.

The thesis adopts a hybrid scope. The synopsis is used as the motivational and planning baseline, but empirical claims are restricted to implemented controls and verified artifacts only. Consequently, the study focuses on crop planning, NPK-aware fertilization, and hierarchical reinforcement learning, while irrigation control, rice-specific localization, and field validation are retained as future work. The latest broad experiment matrix for the Pakistan NPK setup has already been encoded in the repository and standardized through a unified runner with per-run summary JSON outputs, but the full fresh campaign remains pending at the time of this draft build. For that reason, the thesis reports the complete system architecture, methodology, reproducibility design, and historical benchmark context from older audited runs, while explicitly separating those benchmark results from the new NPK campaign that is intended to supply the final comparative chapter.

The main contribution of this work is therefore twofold: first, it transforms the original CyclesGym codebase into a localized, NPK-capable, hierarchically extensible research platform for Pakistan-relevant decision support; second, it defines a defensible thesis-writing and experiment-reporting workflow in which every major claim is anchored either in completed code, completed artifacts, or clearly bounded future work. This structure prevents over-claiming while preserving the scientific value of the implemented engineering advances.
\clearpage
